% Entregable 3 de la Práctica de Laboratorio 02: matriz de hallazgos y ADR-001.
% El contenido del ADR es el de ADR-001-canales-y-despachador-de-avisos.md; si cambia uno, hay que actualizar el otro.
% Compilar con compilar.sh (tabla3.tex primero y este archivo dos veces).
\documentclass[11pt,a4paper]{article}
% Estilo común del informe y del cuestionario de la Práctica de Laboratorio 02 (proyecto OLA).
\usepackage[T1]{fontenc}
\usepackage[utf8]{inputenc}
\usepackage{textcomp}
\DeclareUnicodeCharacter{2212}{\textminus}
\DeclareUnicodeCharacter{2192}{\ensuremath{\rightarrow}}
\DeclareUnicodeCharacter{2264}{\ensuremath{\leq}}
\DeclareUnicodeCharacter{2248}{\ensuremath{\approx}}
\usepackage[spanish,es-tabla,es-nodecimaldot,es-noindentfirst]{babel}
\usepackage[sc]{mathpazo}
\usepackage[scaled=0.92]{helvet}
\renewcommand{\ttdefault}{txtt}
\linespread{1.05}
\usepackage[a4paper,left=2.3cm,right=2.3cm,top=2.7cm,bottom=2.5cm,headsep=0.9cm]{geometry}
\usepackage{microtype}
\usepackage[table]{xcolor}
\usepackage{graphicx}
\usepackage{array,tabularx,longtable,multirow,makecell}
\usepackage{titlesec}
\usepackage{fancyhdr}
\usepackage{lastpage}
\usepackage{enumitem}
\usepackage{listings}
\usepackage[font=small,labelfont={sf,bf},labelsep=colon,skip=6pt]{caption}
\usepackage{pdfpages}
\usepackage{afterpage}
\usepackage[hidelinks,pdfusetitle]{hyperref}
\usepackage{xurl}
\setlength{\emergencystretch}{3em}

% ------------------------------------------------------------------ colores
\definecolor{azul}{HTML}{0F3849}
\definecolor{verdeazul}{HTML}{1F6F80}
\definecolor{fondo}{HTML}{E3EEF0}
\definecolor{fondotenue}{HTML}{F3F7F8}
\definecolor{borde}{HTML}{B9CBCC}
\definecolor{gris}{HTML}{465B62}
\definecolor{rosa}{HTML}{E11D48}
\definecolor{ambar}{HTML}{D97706}
\definecolor{verde}{HTML}{059669}

% ------------------------------------------------------------------ títulos
\titleformat{\section}{\sffamily\bfseries\Large\color{azul}}{\thesection.}{0.6em}{}[{\color{verdeazul}\titlerule[0.6pt]}]
\titleformat{\subsection}{\sffamily\bfseries\large\color{verdeazul}}{\thesubsection}{0.8em}{}
\titleformat{\subsubsection}{\sffamily\bfseries\normalsize\color{azul}}{\thesubsubsection}{0.8em}{}
\titlespacing*{\section}{0pt}{1.6em}{0.9em}
\titlespacing*{\subsection}{0pt}{1.3em}{0.6em}
\setcounter{tocdepth}{2}

% ------------------------------------------------------------------ encabezado y pie
\pagestyle{fancy}
\fancyhf{}
\renewcommand{\headrulewidth}{0.4pt}
\renewcommand{\headrule}{{\color{borde}\hrule width\headwidth height\headrulewidth}}
\fancyhead[L]{\sffamily\footnotesize\color{gris}Calidad de Software · Práctica de Laboratorio 02}
\fancyhead[R]{\sffamily\footnotesize\color{gris}Proyecto OLA}
\fancyfoot[C]{\sffamily\footnotesize\color{gris}Página \thepage\ de \pageref*{LastPage}}
\fancypagestyle{plain}{\fancyhf{}\renewcommand{\headrulewidth}{0pt}%
  \fancyfoot[C]{\sffamily\footnotesize\color{gris}\thepage}}

% ------------------------------------------------------------------ texto
\setlength{\parindent}{0pt}
\setlength{\parskip}{0.55em}
\setlist{itemsep=0.25em,topsep=0.3em}
\newcommand{\ev}[1]{\textsf{\textbf{#1}}}                 % identificador de evidencia
\newcommand{\cod}[1]{\texttt{#1}}                          % código en línea
\newcommand{\negrita}[1]{\textbf{#1}}

% ------------------------------------------------------------------ flotantes
\renewcommand{\topfraction}{0.9}
\renewcommand{\bottomfraction}{0.8}
\renewcommand{\textfraction}{0.07}
\renewcommand{\floatpagefraction}{0.7}
\setcounter{topnumber}{3}\setcounter{bottomnumber}{2}\setcounter{totalnumber}{4}
\makeatletter\setlength{\@fptop}{0pt}\setlength{\@fpbot}{0pt plus 1fil}\setlength{\@fpsep}{16pt}\makeatother
\usepackage[section]{placeins}

% ------------------------------------------------------------------ tablas
\setlength{\arrayrulewidth}{0.4pt}
\arrayrulecolor{borde}
\renewcommand{\arraystretch}{1.25}
\newcolumntype{L}[1]{>{\raggedright\arraybackslash}p{#1}}
\newcolumntype{C}[1]{>{\centering\arraybackslash}p{#1}}
\newcolumntype{Y}{>{\raggedright\arraybackslash}X}
\newcommand{\cab}[1]{\textsf{\bfseries\color{white}#1}}    % celda de cabecera
\newcommand{\filacab}{\rowcolor{azul}}
\newcommand{\filatotal}{\rowcolor{fondo}}
\newcommand{\clave}[1]{\cellcolor{fondo}\textsf{\bfseries\small #1}}

% ------------------------------------------------------------------ fragmentos de código
\lstdefinestyle{ola}{
  language=Python,
  basicstyle=\ttfamily\footnotesize,
  keywordstyle=\color{verdeazul}\bfseries,
  stringstyle=\color{azul},
  commentstyle=\color{gris}\itshape,
  numbers=left, numberstyle=\ttfamily\tiny\color{gris}, numbersep=8pt,
  frame=single, rulecolor=\color{borde}, backgroundcolor=\color{fondotenue},
  xleftmargin=1.6em, framexleftmargin=1.4em, aboveskip=0.8em, belowskip=0.4em,
  columns=fullflexible, keepspaces=true, showstringspaces=false,
  breaklines=true, breakatwhitespace=true, tabsize=4,
  float=htbp, floatplacement=htbp,
  morekeywords={self,None,True,False,Protocol},
}
\lstset{style=ola}
\renewcommand{\lstlistingname}{Fragmento}

% ------------------------------------------------------------------ figuras en página horizontal
% \paginafigura{archivo}{pie}: inserta un PDF A4 horizontal como página propia, con su pie y su número.
\newcommand{\paginafigura}[2]{%
  \includepdf[fitpaper, pagecommand={\thispagestyle{empty}}, picturecommand*={%
    \put(42,13){\parbox{690bp}{\sffamily\footnotesize\color{gris}#2}}%
    \put(800,13){\makebox[0pt][r]{\sffamily\footnotesize\color{gris}Página \thepage\ de \pageref*{LastPage}}}%
  }]{#1}}

% ------------------------------------------------------------------ portada
% \portada{título del documento}{descripción}
\newcommand{\portada}[2]{%
  \hypersetup{pageanchor=false}% la portada no lleva número ni ancla propia
  \begin{titlepage}
    \centering\sffamily\setlength{\parskip}{0pt}
    {\large\bfseries\color{azul}UNIVERSIDAD LA SALLE\par}
    \vspace{0.7cm}
    \includegraphics[width=3.2cm]{figuras/logo_lasalle.png}\par
    \vspace{0.6cm}
    {\large Facultad de Ingenierías y Arquitectura\par}
    \vspace{0.15cm}
    {\large Escuela Profesional de Ingeniería de Software\par}
    \vspace{0.45cm}
    {Curso: Calidad de Software\par}
    \vspace{0.15cm}
    {Docente: Maribel Molina Barriga\par}
    \vspace{1.1cm}
    {\color{verdeazul}\rule{0.62\textwidth}{0.8pt}\par}
    \vspace{0.7cm}
    {\LARGE\bfseries\color{azul}PRÁCTICA DE LABORATORIO 02\par}
    \vspace{0.5cm}
    {\huge\bfseries\color{azul}\begin{minipage}{0.75\textwidth}\centering #1\end{minipage}\par}
    \vspace{0.6cm}
    {\Large\color{verdeazul}\begin{minipage}{0.62\textwidth}\centering Proyecto OLA — Observatorio Litoral de Anomalías térmicas\end{minipage}\par}
    \vspace{0.7cm}
    {\color{verdeazul}\rule{0.62\textwidth}{0.8pt}\par}
    \vspace{0.6cm}
    {\rmfamily\itshape\color{gris}\begin{minipage}{0.66\textwidth}\centering #2\end{minipage}\par}
    \vspace{0.9cm}
    {\bfseries Integrantes\par}
    \vspace{0.35cm}
    {\renewcommand{\baselinestretch}{1.25}\selectfont
    Frederick Mares Graos\par
    Jhordan Huamani Huamani\par
    Jorge Ortiz Castañeda\par
    Miguel Angel Flores Leon\par
    Piero Adrian Delgado Chipana\par}
    \vfill
    Arequipa, octubre de 2026\par
  \end{titlepage}\hypersetup{pageanchor=true}}

\hypersetup{pdftitle={Práctica de Laboratorio 02 — Matriz de hallazgos y ADR},
            pdfauthor={Equipo OLA}, pdfsubject={Calidad de Software}}

\begin{document}

\portada{Matriz de hallazgos y ADR}
        {Matriz de hallazgos priorizada del módulo «Alertas y avisos» y registro de decisión de arquitectura ADR-001}

{\hypersetup{linkcolor=black}\renewcommand{\contentsname}{\sffamily\color{azul}Índice}\tableofcontents}
\thispagestyle{plain}
\newpage

% =====================================================================
\section{Matriz de hallazgos}

La revisión del módulo «Alertas y avisos» (commit \cod{d5c054f}) produjo 12 hallazgos, que la matriz de la página siguiente ordena por prioridad (\ev{E-04}). De ellos, 5 quedan resueltos por el parche de refactorización (H-01 a H-05), 5 siguen pendientes (H-06, H-07, H-09, H-10 y H-12) y 2 se aceptaron con su razón (H-08 y H-11). La misma matriz, con sus fórmulas, está en la hoja «Hallazgos» de \cod{Matriz\_Hallazgos\_OLA.xlsx}.

\negrita{Escala.} Severidad: Alta = impacto 4--5; Media = 3; Baja = 1--2. Impacto: 1 = cosmético a 5 = bloquea cambios previsibles o impide probar. Esfuerzo: 1 = menos de una hora a 5 = varios días. Prioridad: orden por impacto $\times$ 10 $-$ esfuerzo.

\negrita{Nota sobre la numeración.} La matriz corresponde a la Tabla 3 del informe técnico; conserva su título original.

\afterpage{\paginafigura{tabla3.pdf}{}}

% =====================================================================
\clearpage
\section{ADR-001: Separar el envío de avisos en canales, almacén y despachador}

\begin{center}
\small
\rowcolors{1}{white}{fondotenue}
\begin{tabularx}{\textwidth}{L{4.2cm}Y}
\hline
\clave{Estado} & Aceptado (pendiente de aplicar el parche en una rama)\\ \hline
\clave{Fecha de la propuesta} & 2026-10-05\\ \hline
\clave{Fecha de aceptación} & 2026-10-06\\ \hline
\clave{Proyecto} & OLA — Observatorio Litoral de Anomalías térmicas\\ \hline
\clave{Alcance} & Backend, módulo «Alertas y avisos» (RF-01, RF-03)\\ \hline
\clave{Versión evaluada} & commit \cod{d5c054f} (2026-09-13)\\ \hline
\clave{Decisores} & Equipo OLA; aprobó Miguel Angel Flores Leon (Product Owner)\\ \hline
\clave{Hallazgos que resuelve} & H-01, H-02, H-03, H-04, H-05\\ \hline
\end{tabularx}
\end{center}

\subsection{Contexto}

\cod{services/notification\_service.py} reúne cinco responsabilidades: decide qué avisos corresponden a cada episodio, los inserta en la base, redacta el mensaje, lo envía por correo y marca los avisos como leídos. La revisión de diseño encontró cuatro problemas que nacen de ahí:

\begin{enumerate}
  \item \negrita{Los canales no son extensibles.} \cod{CANALES} es una tupla fija, \cod{send\_pending} solo sabe enviar correo y \cod{notifications\_repo} filtra por \cod{EMAIL} en dos consultas. Añadir un canal obliga a modificar el servicio, el repositorio, el router y las dependencias (H-02).
  \item \negrita{El servicio construye SQL.} Contiene 4 sentencias \cod{select} / \cod{insert … on\_conflict} y repite dos veces el mismo bloque de inserción (H-01, H-03).
  \item \negrita{La lógica de envío no se puede probar sin PostgreSQL.} Los reintentos, el registro del fallo y el límite de la tanda tienen 0 pruebas unitarias y 34 de integración, que tardan entre 12.5 y 12.9~s (H-05).
  \item \negrita{La abstracción \cod{Clock} existe pero no se usa.} Ningún módulo de producción la importa; el servicio llama a \cod{datetime.now(UTC)} en 3 lugares (H-04).
\end{enumerate}

La SRS retiró el SMS de RF-03 por su costo por mensaje, de modo que hoy solo existen el correo y el aviso dentro de la aplicación. El equipo considera que el cambio más previsible del módulo es añadir otro canal de aviso.

Restricciones: el equipo no puede ejecutar una reescritura grande antes del Hito~2, y las 442 pruebas existentes del backend deben seguir pasando sin modificarlas.

\subsection{Opciones consideradas}

\begin{center}
\small
\rowcolors{2}{white}{fondotenue}
\begin{tabularx}{\textwidth}{L{4.0cm}Y}
\filacab \cab{Opción} & \cab{Descripción}\\
\negrita{A. Puertos y canales} & Un contrato \cod{Channel} por canal y un registro de canales; un contrato \cod{NotificationStore} para la persistencia; un \cod{NotificationDispatcher} que recibe almacén, canales y reloj.\\
\negrita{B. Cambio mínimo} & Mover el SQL al repositorio y añadir un condicional por canal dentro de \cod{send\_pending}.\\
\negrita{C. Cola asíncrona} & Publicar cada aviso en una cola y enviarlo con un trabajador en segundo plano.\\
\end{tabularx}
\end{center}

Matriz ponderada (1 = peor, 5 = mejor; en esfuerzo y riesgo, 5 es menor esfuerzo y menor riesgo):

\begin{center}
\small
\rowcolors{2}{white}{fondotenue}
\begin{tabular}{L{6.0cm}C{1.4cm}C{1.6cm}C{1.6cm}C{1.6cm}}
\filacab \cab{Criterio} & \cab{Peso} & \cab{A} & \cab{B} & \cab{C}\\
Modificabilidad & 30 \% & 5 & 3 & 5\\
Comprobabilidad & 25 \% & 5 & 2 & 4\\
Esfuerzo de implementación & 20 \% & 3 & 5 & 1\\
Riesgo de regresión & 15 \% & 4 & 5 & 2\\
Comprensibilidad para el equipo & 10 \% & 3 & 4 & 2\\
\filatotal\textsf{\textbf{Total ponderado}} & & \textsf{\textbf{4.25}} & 3.55 & 3.20\\
\end{tabular}
\end{center}

\subsection{Decisión}

Se adopta la \negrita{opción A}.

\begin{itemize}
  \item \cod{domain/notices.py} define los conceptos del aviso sin dependencias de infraestructura: los enumerados de canal, tipo y estado (que antes estaban en \cod{db/models.py}), \cod{plan\_rows} y \cod{compose}.
  \item \cod{services/channels.py} define el contrato \cod{Channel}, los canales \cod{EmailChannel} e \cod{InAppChannel}, el \cod{ChannelRegistry} y la función \cod{default\_channels}, único punto donde se decide qué canales tiene el sistema.
  \item \cod{services/notification\_dispatcher.py} contiene \cod{NotificationDispatcher}, que depende solo de los protocolos \cod{NotificationStore}, \cod{Channel} y \cod{Clock}.
  \item \cod{repositories/notifications\_repo.py} concentra la inserción (\cod{insert\_pending}) y aporta \cod{SqlNotificationStore}, la implementación del almacén sobre PostgreSQL.
  \item \cod{services/notification\_service.py} queda como coordinador: fija el límite de cada transacción y conserva sus firmas públicas (\cod{send\_pending} sigue existiendo y delega en \cod{dispatch\_pending}), de modo que las pruebas existentes no cambian.
  \item \cod{api/deps.py} añade \cod{get\_channels}, y \cod{api/routers/alerts.py} recibe los canales por inyección (\cod{ChannelsDep}) en lugar del \cod{Mailer}. Es el único router que cambia, en 3 líneas.
\end{itemize}

\subsection{Consecuencias}

\subsubsection*{Positivas}

\begin{itemize}
  \item \negrita{Modificabilidad.} Se añadió un canal de prueba (una pasarela por \emph{webhook}) con la misma especificación en ambos diseños. Los módulos existentes modificados bajan de 6 a 4 y las líneas de código existente reescritas, de 13 a 4. El servicio, el despachador, el repositorio y el router no cambian. La complejidad del bucle de envío se mantiene en 4, mientras que en el diseño actual sube a 6.
  \item \negrita{Comprobabilidad.} La lógica de envío pasa de 0 a 12 pruebas unitarias; con las de canales y avisos son 30, que se ejecutan en 0.04~s sin base de datos. \cod{domain/notices.py} se importa sin ninguna biblioteca externa, y el despachador y los canales, sin SQLAlchemy, FastAPI ni psycopg.
  \item \negrita{Cohesión.} \cod{notification\_service} baja de 141 a 73 líneas y su complejidad acumulada, de 20 a 12; deja de construir SQL (4 sentencias a 0).
  \item \cod{Clock} pasa a usarse (Ca: 0 a 2).
\end{itemize}

\subsubsection*{Negativas}

\begin{itemize}
  \item \negrita{Más código.} El módulo revisado crece de 1\,011 a 1\,173 líneas (+16 \%) y aparecen 3 módulos nuevos.
  \item \negrita{Más indirección.} Seguir un envío exige leer cuatro archivos en lugar de uno. \cod{notification\_service} depende ahora de 8 módulos internos en lugar de 5.
  \item \negrita{No reduce las pruebas que hay que actualizar al añadir un canal.} En el escenario del \emph{webhook} dejan de pasar las mismas 11 pruebas en ambos diseños, porque cuentan los avisos por episodio o el resumen del envío, que cambian por especificación. El diseño propuesto suma una más, \cod{test\_solo\_el\_correo\_requiere\_envio}, que fija qué canales requieren envío.
  \item \negrita{Los canales siguen cargando la configuración completa.} \cod{channels} importa \cod{mail}, y \cod{mail} importa \cod{Settings}; por eso el despachador y los canales cargan pydantic y smtplib (H-07).
  \item \negrita{El enumerado de canal sigue en PostgreSQL.} Añadir un canal exige una migración en cualquier diseño (H-11).
  \item Quien incorpore un canal debe registrarlo en \cod{default\_\allowbreak channels}; si lo olvida, sus avisos quedan pendientes. La prueba \cod{test\_\allowbreak los\_\allowbreak canales\_\allowbreak por\_\allowbreak defecto\_\allowbreak cubren\_\allowbreak todos\_\allowbreak los\_\allowbreak que\_\allowbreak se\_\allowbreak guardan} falla en ese caso.
\end{itemize}

\subsubsection*{Pendiente}

\begin{itemize}
  \item Aplicar \cod{refactorizacion\_avisos.patch} en una rama del repositorio y fusionarla.
  \item Aplicar el mismo criterio a \cod{alert\_service}, que también construye SQL (H-06).
  \item Fijar el límite de las transacciones en un solo nivel (H-12): \cod{notify\_open\_alerts\_of} sigue sin confirmar y lo hace el router.
  \item \cod{count\_by\_status} sigue contando solo el correo.
\end{itemize}

\subsection{Verificación}

\begin{center}
\small
\rowcolors{2}{white}{fondotenue}
\begin{tabularx}{\textwidth}{L{6.4cm}Y}
\filacab \cab{Comprobación} & \cab{Resultado}\\
Pruebas unitarias previas & 212 de 212 pasan\\
Pruebas de integración previas & 230 de 230 pasan, con el dataset montado en \cod{/dataset}\\
Pruebas unitarias nuevas & 30 de 30 pasan en 0.04~s\\
\cod{ruff check .} con la configuración del proyecto & Sin observaciones, antes y después\\
\cod{mypy src} (estricto, con el complemento de pydantic) & 0 errores antes (61 archivos) y después (64 archivos)\\
Cobertura de sentencias del backend & 95.11 \% antes y 95.39 \% después\\
Cobertura de los módulos del alcance & 95.63 \% antes y 96.32 \% después\\
El parche se aplica sobre \cod{d5c054f} & \cod{git apply -{}-check} sin conflictos; 472 pruebas pasan\\
\end{tabularx}
\end{center}

Entorno de la verificación: imagen \cod{ola-api} del proyecto (Python 3.13.15, pytest 9.1.1, ruff 0.16.6, mypy 2.3.1, bcrypt 5.0.0) contra PostgreSQL 17.11, con la configuración de \cod{.env.example} y secretos propios. La base de datos corrió en un contenedor desechable en lugar del servicio \cod{db} de \cod{compose.yml}.

\subsection{Referencias}

\begin{itemize}
  \item Informe de la Práctica de Laboratorio 02, secciones 3 a 6 (evidencias \ev{E-01} a \ev{E-11}).
  \item \cod{refactorizacion\_avisos.patch}: el cambio completo.
  \item \cod{escenario\_webhook/}: el escenario de cambio aplicado a los dos diseños.
  \item \cod{evidencias/}: salidas en bruto de pytest, ruff, mypy, la cobertura y las métricas.
  \item M. Nygard, «Documenting Architecture Decisions», 2011.
  \item R. C. Martin, \emph{Clean Architecture}, Pearson, 2017, caps. 8 (OCP) y 11 (DIP).
  \item ISO/IEC 25010:2023, característica de mantenibilidad.
\end{itemize}

\end{document}
